% Part: proof-theory
% Chapter: sequent-calculus
% Section: admissible-derivable

\documentclass[../../../include/open-logic-section]{subfiles}

\begin{document}

\olfileid{pt}{seq}{adm}

\olsection{مجازې او د اشتقاق وړ قاعدې}

د ثبوت د تيورۍ ډېرې پايلې د دې پوښتنې په اړه دي، يا د هغې په اړه
پايلې کاروي، چې آيا هر څۀ چې د يوې ځانګړې قاعدې په کارولو سره
!!{prove}d کېدے شي، له هغې قاعدې پرته هم !!{prove}d کېدے شي۔
د پرېکون د له منځه وړلو قضيه ئې تر ټولو پېژندلے مثال دے۔ دا ښيي
چې هر سېکوېنټ چې د يو سېکوېنټي نظام د قواعدو او~\Cut{} په
کارولو سره !!{prove}d کېدے شي، له~\Cut پرته هم !!{prove}d کېدے شي۔
ځکه چې دا ډول خاصيت ډېر مهم دے، خپل نوم لري۔

\begin{defn}
قاعده~$R$ په يو سېکوېنټي نظام کښې \emph{مجاز} ده، که هر هغه
سېکوېنټ چې په دغه نظام کښې له قاعدې~$R$ سره ثابتېدے شي،
له~$R$ پرته هم ثابتېدے شي۔
\end{defn}

\begin{ex}\ollabel{ex:land-deriv}
قاعده $\LeftR{\land}$ په \Log{G1c} او
\Log{G3c} کښې بېلابېلې بڼې لري۔ په \Log{G1c} کښې د دې قاعدې
دوه بڼې دي: په يوه کښې د او تړنې کيڼ غړے $!A$ د اصلي
!!{formula}~$!A \land !B$ په مقدمه کښې دے، او په بله کښې ئې ښے
د او تړنې غړے~$!B$ دے:
\[
\Axiom$ !A, \Gamma \fCenter \Delta$
\RightLabel{$\LeftR{\land}_1$}
\UnaryInf$ !A \land !B, \Gamma \fCenter \Delta$
\DisplayProof
\qquad
\Axiom$!B, \Gamma \fCenter \Delta$
\RightLabel{$\LeftR{\land}_2$}
\UnaryInf$!A \land !B, \Gamma \fCenter \Delta$
\DisplayProof
\]
په مقابل کښې، \Log{G3c} يوازې يوه قاعده لري:
\[
\Axiom$ !A, !B, \Gamma \fCenter \Delta$
\RightLabel{$\LeftR{\land}_3$}
\UnaryInf$ !A \land !B, \Gamma \fCenter \Delta$
\DisplayProof
\]
اوس دا پوښتنه کولے شو: «آيا د \Log{G3c} قاعده
$\LeftR{\land}_3$ په \Log{G1c} کښې مجاز ده؟» هو، ده:
هر هغه استنتاج چې $\LeftR{\land}_3$ کاروي، يوازې د~\Log{G1c}
د قواعدو په کارولو سره نېغ په نېغه تمثيلولے شو۔ د
$\LeftR{\land}_1$ او $\LeftR{\land}_2$ ترڅنګ به جوړښتي
قاعدې~$\LeftR{\Contraction}$ ته هم اړتيا ولرو:
\[
\Axiom$!A, !B, \Gamma \fCenter \Delta$
\RightLabel{$\LeftR{\land}_1$}
\UnaryInf$!A \land !B, !B, \Gamma \fCenter \Delta$
\RightLabel{$\LeftR{\land}_2$}
\UnaryInf$!A \land !B, !A \land !B, \Gamma \fCenter \Delta$
\RightLabel{\LeftR{\Contraction}}
\UnaryInf$!A \land !B, \Gamma \fCenter \Delta$
\DisplayProof
\]
\textbf{د سرچينې سمون (OLCMP-222):} په وروستۍ نتيجه کښې
د اصلي او تړنې او پاتې سياق ترمنځ ورکه کامه ورزياته شوه؛
د انقباض قاعده او نور ټول فورمولونه عين دي۔
\end{ex}

د يوې قاعدې په کارولو سره شوي استنتاجونه د نورو قواعدو په وسيله
تمثيلول د دې ښودلو يوه لاره ده چې قاعده مجاز ده۔ خو دا يوازينۍ
لاره نۀ ده؛ ځينې قاعدې مجازې دي، خو داسې نۀ شي تمثيلېدې۔
هغه قاعدې چې داسې تمثيلېدې شي، \emph{د اشتقاق وړ} بلل کېږي۔

\begin{defn}
قاعده~$R$،
\[
\AxiomC{$S_1$}
\AxiomC{$\dots$}
\AxiomC{$S_n$}
\RightLabel{$R$}
\TrinaryInfC{$S$}
\DisplayProof
\]
په يو نظام کښې \emph{د اشتقاق وړ} بلل کېږي، که د سېکوېنټ~$S$
يو سکيماتي !!{proof} موجود وي چې د نظام قواعد او بديهيات
او د~$S_1$، \dots،~$S_n$ د بڼو اضافي ابتدايي سېکوېنټونه کاروي۔
\end{defn}

په \olref{ex:land-deriv} کښې ورکړے شوے !!{proof} ښيي چې
$\LeftR{\land}_3$ په~\Log{G1c} کښې د اشتقاق وړ قاعده ده؛ ځکه
دا !!{proof} د $\LeftR{\land}_3$ د پايلې يو سکيماتي !!{proof} دے
چې د $\LeftR{\land}_3$ له مقدماتو څخه، د ابتدايي سېکوېنټونو په
توګه، يوازې د \Log{G1c} د قواعدو په کارولو جوړ شوے دے۔
(د \Log{G1c} هېڅ بديهيه ئې نۀ ده کارولې۔)

\begin{prop}\ollabel{prop:lor-G3-adm}
د~\Log{G3c} قاعدې $\LeftR{\land}$ او $\RightR{\lor}$
په~\Log{G1c} کښې د اشتقاق وړ دي۔
\end{prop}

\begin{proof}
  د $\LeftR{\land}$ ثبوت په \olref{ex:land-deriv} کښې ورکړے شوے دے۔
  د $\RightR{\lor}$ ثبوت د تمرين په توګه پرېښودل شوے دے۔
\end{proof}

\begin{prob}
وښيئ چې د \Log{G3c} قاعده $\RightR{\lor}$
په~\Log{G1c} کښې د اشتقاق وړ ده۔
\end{prob}

\begin{prob}
فرض کړئ $\liff$ يو تعريف شوے !!{operator} دے: $!A \liff !B$
د $(!A \lif !B) \land (!B \lif !A)$ لنډه ليکنه ده۔ وښيئ چې
لاندې قاعدې
\begin{enumerate}
\item \Axiom$\Gamma, !A, !B \fCenter \Delta$
\Axiom$\Gamma \fCenter \Delta, !A, !B$
\RightLabel{\LeftR{\liff}}
\BinaryInf$!A \liff !B, \Gamma \fCenter \Delta$
\DisplayProof
\textbf{د سرچينې سمون (OLCMP-223):} د کيڼې دوه‌شرطي
قاعدې اصلي پايله په غلطۍ ښي اړخ کښې وه؛ همدغه
دوه‌شرطي فورمول کيڼ اړخ ته سم شوے۔ دواړه مقدمې،
د ښۍ قاعدې صحيح پايله او د دواړو نظامونو تمرين عين دي۔
\item
\Axiom$!A, \Gamma \fCenter \Delta, !B$
\Axiom$!B, \Gamma \fCenter \Delta, !A$
\RightLabel{\RightR{\liff}}
\BinaryInf$\Gamma \fCenter \Delta, !A \liff !B$
\DisplayProof
\end{enumerate}
په (a)~\Log{G1c} او (b)~\Log{G3c} کښې د اشتقاق وړ دي۔
\end{prob}

\begin{ex}\ollabel{ex:weak-adm}
د \Log{G1c} قاعده $\RightR{\Weakening}$،
\[
\Axiom$ \Gamma \fCenter \Delta$
\RightLabel{\RightR{\Weakening}}
\UnaryInf$ \Gamma \fCenter \Delta, !A$
\DisplayProof
\]
په~\Log{G3c} کښې مجاز ده۔ د ثبوت نظر ساده دے: فرض کړئ يو
!!a{proof}~$\pi$ د مقدمې $\Gamma \Sequent \Delta$ دپاره په~\Log{G3c}
کښې موجود دے۔ \textbf{د سرچينې سپيناوی (OLCMP-224):} لومړے د ټول
ثبوت ځانګړي متغيرونه داسې تازه نومونو ته واړوئ چې
په ورزياتېدونکي فورمول کښې نۀ راځي۔ بيا سياق پراخ
کړئ؛ ګنې د کمي قواعدو د ځانګړي متغير شرط ماتېدلے شي۔
دا نوم بدلول د ثبوت لوړوالے نۀ بدلوي۔
!!{formula}~$!A$ په~$\pi$ کښې د هر سېکوېنټ ښۍ
خوا ته ورزيات کړئ۔
بديهيات بديهيات پاتې کېږي، او سم استنتاجونه لا هم سم استنتاجونه
دي۔ د نوي !!{proof} وروستنے سېکوېنټ، چې اصل ئې وروستۍ !!{formula} بولي، دا دے: $\Gamma \Sequent
\Delta, !A$۔

خو قاعده $\RightR{\Weakening}$ په~\Log{G3c} کښې د اشتقاق وړ
نۀ ده۔ ځکه فرض کړئ چې د اشتقاق وړ وه؛ يعنې فرض کړئ د
$\Gamma \Sequent
\Delta, !A$ يو !!a{proof} موندل کېدے شو چې يوازې د~\Log{G3c}
بديهيات او قواعد او، ښايي، اضافي ابتدايي سېکوېنټ~$\Gamma \Sequent \Delta$
کاروي۔ که دا !!{proof} په بشپړه توګه سکيماتي وي، لکه د
$\RightR{\Weakening}$ د اشتقاق وړ کېدو دپاره چې اړين دي، نو دا په
يو !!a{proof} بدلولے شو چې د~$\Sequent !A$ دپاره وي، که
$\Gamma$ او~$\Delta$ لرې کړو۔ په دې سره اضافي ابتدايي
سېکوېنټونه~$\Gamma \Sequent \Delta$ په تش
سېکوېنټ~$\quad \Sequent \quad$ بدلېږي۔ ځکه چې تش سېکوېنټ هېڅ
!!{formula}s نۀ لري، خو د~\Log{G3c} ټول قواعد په خپلو مقدماتو کښې
لږ تر لږه يوه فعاله !!{formula} غواړي، نو تش سېکوېنټ په دې سکيماتي
!!{proof} کښې د هېڅ استنتاج مقدمه نۀ شي کېدے۔ په نتيجه کښې،
!!{proof} يوازې هغه وخت سکيماتي کېدے شي چې په رښتيا اضافي
سېکوېنټ $\Gamma \Sequent \Delta$ د ابتدايي سېکوېنټ په توګه نۀ کاروي۔
په بله وينا، دا به د $\Gamma \Sequent
\Delta, !A$ يو سکيماتي اشتقاق وي چې يوازې د~\Log{G3c} بديهيات
او قواعد کاروي؛ او موږ به ښودلې وي چې د دې بڼې هر سېکوېنټ
په~\Log{G3c} کښې يو !!a{proof} لري۔ خو دا ناشونې ده، ځکه \Log{G3c}
سم دے او نو يوازې هغه سېکوېنټونه !!{prove} کولے شي چې په هر
!!{structure} کښې رښتيني وي۔ که $\Gamma = \Delta
= \emptyset$ او $!A \ident \lfalse$ واخلو، \Log{G3c} به د بېلګې
په توګه سېکوېنټ $\quad \Sequent \lfalse$ !!{prove} کولے شي، حال دا
چې داسې نۀ شي کولے۔
\end{ex}

راځئ د دې پايلې يو ښکلے استقرايي ثبوت ورکړو چې کمزوري کول
په~\Log{G3c} کښې يوه مجازه قاعده ده۔ په حقيقت کښې، لکه شهودي دليل
چې ښيي، يو څۀ قوي پايله هم سمه ده: که $!A$ د~$\pi$ د هر
سېکوېنټ ښۍ خوا ته ورزيات کړو، نو د !!{proof} جوړښت نۀ بدلوو؛
په ځانګړې توګه، د ثبوت لوړوالے نۀ زياتوو۔ ځکه چې دا مهمه ده،
خپل تعريف غواړي۔
\textbf{د ځانګړي متغير شرط (OLCMP-224):} د سياقونو له
پراخولو مخکې ځانګړي متغيرونه له ټولو ورزياتېدونکو
فورمولونو څخه تازه کړئ؛ کيڼ او ښي کمزوري کول د دې
نوم بدلولو وروسته د ثبوت لوړوالے ساتي۔

\begin{defn}
قاعده~$R$،
\[
\AxiomC{$S_1$}
\AxiomC{$\dots$}
\AxiomC{$S_n$}
\RightLabel{$R$}
\TrinaryInfC{$S$}
\DisplayProof
\]
په يو نظام کښې \emph{!!{height} ساتونکې مجازه} (يا
\emph{!!{hp}-مجازه}) ده، که هر کله $\Proves[h_i] S_i$ د
$i = 1$، \dots،~$n$ دپاره وي، نو $\Proves_m S$ هم وي، چې
$m = \max(h_1, \dots,
h_n)$ دے۔
\end{defn}

په ياد ولرئ چې د يوې قاعدې د !!{height} ساتونکې مجازې کېدو دپاره
دومره بس دي چې، هر کله مقدمې~$S_i$ !!{proof}s~$\pi_i$ ولري
چې !!{height}~$h_i = \pheight{\pi_i}$ لري،
د پايلې يو !!a{proof}~$\pi$ موجود وي چې !!{height}~$\pheight{\pi}$
ئې د~$h_i$ له اعظمي ارزښت څخه زيات نۀ وي۔ په ځانګړې توګه،
که يوازې يوه مقدمه~$S_1$ وي، نو دومره بس دي چې
$\pheight{\pi} \le \pheight{\pi_1}$ وي۔ د $\RightR{\Weakening}$
او $\LeftR{\Weakening}$ په صورت کښې په رښتيا
$\pheight{\pi} = \pheight{\pi_1}$ لرو، خو دا اړينه نۀ ده۔

\begin{prop}\ollabel{prop:weak-G3c-adm}
د \Log{G1c} قاعدې $\RightR{\Weakening}$ او $\LeftR{\Weakening}$
په~\Log{G3c} کښې !!{height} ساتونکې مجازې دي۔
\end{prop}

\begin{proof}
  بايد وښيو چې که $\Log{G3c} \Proves \Gamma \Sequent \Delta$
  له يو !!{proof}~$\pi$ سره وي چې !!{height}~$\pheight{\pi} = n$ لري،
  نو يو \Log{G3c}-!!{proof} $\pi'$ د $\Gamma \Sequent \Delta, !A$ دپاره
  موجود دے چې $\pheight{\pi'} = n$ وي۔ دا پر~$n$ د استقرا په
  وسيله ثابتوو۔ که $n=0$ وي، نو $\pi$ يوازې بديهيه
  $\Gamma \Sequent \Delta$ ده (يعنې يو اتومي $!C$ په دواړو
  $\Gamma$ او~$\Delta$ کښې راځي، يا $\lfalse$ په $\Gamma$ کښې راځي)، او $\Gamma
  \Sequent \Delta, !A$ هم بديهيه ده۔ که $n = \pheight{\pi} >0$
  وي، نو !!{proof}~$\pi$ يو وروستے استنتاج لري۔ د کارول شوې
  قاعدې له مخې بېلابېل صورتونه جلا کوو۔ يوازې يوه بېلګه ورکوو:
  فرض کړئ وروستے استنتاج $\RightR{\land}$ و۔ نو
  $\Delta = \Delta', !B \land !C$ دے، او د وروستي استنتاج
  نېغې فرعي !!{proof}s $\pi_1$ او~$\pi_2$ په ترتيب سره وروستي
  سېکوېنټونه $\Gamma \Sequent \Delta', !B$ او $\Gamma \Sequent
  \Delta', !C$ لري۔ په بله وينا، $\pi$ داسې ښکاري:
  \[
  \AxiomC{}
  \RightLabel{$\pi_1$}
  \Deduce$\Gamma \fCenter \Delta', !B$
  \AxiomC{}
  \RightLabel{$\pi_2$}
  \Deduce$\Gamma \fCenter \Delta', !C$
  \RightLabel{\RightR{\land}}
  \BinaryInf$\Gamma \fCenter \Delta', !B \land !C$
  \DisplayProof
  \]
  همداراز $n =
  \max(\pheight{\pi_1}, \pheight{\pi_2}) + 1$ لرو۔ نو $\pheight{\pi_1} <
  n$ او $\pheight{\pi_2} < n$ دي، او د استقرا فرضيه تطبيقېږي:
  !!{proof}s $\pi_1'$ او~$\pi_2'$ د $\Gamma \Sequent
  \Delta', !B, !A$ او $\Gamma \Sequent \Delta', !C, !A$ دپاره موجود دي۔
  د $\RightR{\land}$ قاعدې په کارولو سره يو !!a{proof}~$\pi'$ د
  $\Gamma
  \Sequent \Delta', !B \land !C, !A$ دپاره ترلاسه کوو:
  \[
  \AxiomC{}
  \RightLabel{$\pi_1'$}
  \Deduce$\Gamma \fCenter \Delta', !B, !A$
  \AxiomC{}
  \RightLabel{$\pi_2'$}
  \Deduce$\Gamma \fCenter \Delta', !C, !A$
  \RightLabel{\RightR{\land}}
  \BinaryInf$\Gamma \fCenter \Delta', !B \land !C, !A$
  \DisplayProof
  \]
  موږ $\pheight{\pi'} =
  \max(\pheight{\pi_1'}, \pheight{\pi_2'}) + 1 = \max(\pheight{\pi_1},
  \pheight{\pi_2}) + 1 = \pheight{\pi}$ لرو۔
\end{proof}
\textbf{د ځانګړي متغير شرط (OLCMP-224):} د سياقونو له
پراخولو مخکې ځانګړي متغيرونه له ټولو ورزياتېدونکو
فورمولونو څخه تازه کړئ؛ کيڼ او ښي کمزوري کول د دې
نوم بدلولو وروسته د ثبوت لوړوالے ساتي۔
\textbf{د سرچينې سمون (OLCMP-225):} د صفر لوړوالي
د بديهي اصل په حالت کښې د تشې نښې کيڼ بديهي اصل هم
شامل دے؛ اصل يوازې د ګډ اتومي فورمول حالت يادولو۔
د دواړو بديهي اصولو په سياق کښې کمزوري کول لوړوالے
صفر ساتي۔

\olref{prop:weak-G3c-adm} ډېر ګټور دے۔ د دې د تکراري تطبيق دپاره
يوه ليکنښه معرفي کوو: که $\pi$ د~$\Gamma \Sequent \Delta$
يو !!a{proof} وي، نو $\pi[\Pi \Sequent \Lambda]$ هغه پايله ده
چې د~$\Pi$ له ټولو !!{formula}s سره په کيڼ اړخ کښې د کمزوري کولو
مجازتيا او د~$\Lambda$ له ټولو !!{formula}s سره په ښي اړخ کښې
د کمزوري کولو مجازتيا تطبيق کړو۔ دا د $\Gamma, \Pi
\Sequent \Delta, \Lambda$ يو !!a{proof} دے۔
\textbf{د ځانګړي متغير شرط (OLCMP-224):} د سياقونو له
پراخولو مخکې ځانګړي متغيرونه له ټولو ورزياتېدونکو
فورمولونو څخه تازه کړئ؛ کيڼ او ښي کمزوري کول د دې
نوم بدلولو وروسته د ثبوت لوړوالے ساتي۔

\end{document}
